% ============================================================
% Come si usano, richiami di matematica
% ============================================================
\argomento[R]{Richiami di matematica}{Formule a p.~\pageref*{f:sommatorie}}

\scheda{sommatorie}{
  \usoquando{Nei calcoli di medie e varianze: le costanti escono dalla somma e le somme si spezzano.}
  \usocome{Il cambio di indice serve a riconoscere una somma nota dopo aver semplificato: per esempio la somma di
    tutti i valori di una pmf, che vale 1.}
  \usoesempio{$\sum_x (a x + b)\, p_X(x) = a \sum_x x\, p_X(x) + b \sum_x p_X(x) = a\,E[X] + b$.}
}

\scheda{gauss}{
  \usoquando{Somme di interi consecutivi; media dell'uniforme su $\{0, \ldots, M - 1\}$ (\vv{uniforme}).}
  \usoesempio{$1 + 2 + \cdots + 100 = \frac{100 \cdot 101}{2} = 5\,050$.}
}

\scheda{quadrati}{
  \usoquando{Valore quadratico medio, e quindi varianza, dell'uniforme discreta.}
  \usoesempio{$1 + 4 + 9 = 14 = \frac{3 \cdot 4 \cdot 7}{6}$.}
}

\scheda{geometrica}{
  \usoquando{Somme infinite di probabilità che si riducono di un fattore costante a ogni passo, come «la prima testa
    esce al lancio $n$».}
  \usoattenzione{Vale solo per $\abs{q} < 1$, e bisogna guardare da quale indice parte la somma: da 0 vale
    $\frac{1}{1 - q}$, da 1 vale $\frac{q}{1 - q}$.}
  \usoesempio{$\sum_{k=0}^{\infty} (1/3)^k = \frac{1}{1 - 1/3} = \frac{3}{2}$; \
    $\sum_{n=1}^{\infty} (1/2)^n = \frac{1/2}{1 - 1/2} = 1$.}
}

\scheda{nepero}{
  \usoquando{Con la variabile di Poisson: la serie esponenziale dimostra che la pmf somma a 1 e ne dà la media; il
    limite notevole porta dalla binomiale alla Poisson (\vv{poisson}).}
  \usoesempio{$\sum_{k=0}^{\infty} \frac{\lambda^k}{k!}\, e^{-\lambda} = e^{-\lambda}\, e^{\lambda} = 1$. Con
    $n = 10^6$: $(1 + 1/n)^n \simeq 2{,}71828$.}
}

\scheda{limiti}{
  \usoquando{Per passare dalle frequenze alle probabilità: un'uguaglianza fra frequenze, vera per ogni $n$, diventa
    un'uguaglianza fra probabilità. E nel limite della binomiale.}
  \usoesempio{$f_n(\compl{A}) = 1 - f_n(A)$ per ogni $n$, quindi al limite $P(\compl{A}) = 1 - P(A)$.}
}
